
\chapter{Proposed Asynchronous Flash ADC Design}
\section{Architectural Evolution and Methodology}
The design of the voltage comparator responsible for generating the digital
output code of the analog-to-digital converter was not arrived at directly.
It evolved through a sequence of candidate topologies, each evaluated against
the constraints imposed by the target supply voltage, the required conversion
accuracy, the effective input dynamic range, and the overall demands on circuit
compactness. The following sections document this progression, presenting
each topology that was considered and explaining the physical reasoning that
either led to its adoption or prompted its replacement by a subsequent
candidate.


\section{Phase I: The Baseline Current-Mirror OTA Flash Configuration}
\subsection{Operating Principle and Core Topology}
\subsection{Current Comparator Topology Exploration and Characterization}
The first topology investigated was the current-mirror
flash ADC (CMF-ADC) introduced by Lee~\textit{et al.}~\cite{lee2017},
originally proposed as the sensing front end of a coarse-fine dual-loop digital
low-dropout regulator. Although the application context of that work differs
substantially from the present design, the operating principle of the CMF-ADC
constitutes a direct voltage-to-thermometer-code conversion mechanism that is
directly applicable in the context of multi-bit comparison, making it a natural
starting point for the architectural exploration.

The circuit is divided into two cascaded functional stages: a transconductance
cell and a current-to-code converter. The transconductance cell is implemented
as a differential pair. Its function is to
translate the voltage difference between the input voltage, $V_\mathrm{in}$,
and an adjustable reference voltage, $V_\mathrm{ref,c}$, into an error
current $I_\mathrm{err}$, yielding branch currents
$I_1 = I_B - I_\mathrm{err}$ and $I_2 = I_B + I_\mathrm{err}$, where a
positive $I_\mathrm{err}$ corresponds to $V_\mathrm{in}$ exceeding
$V_\mathrm{ref,c}$. A third current, $I_3$, replicates $I_1$ through a
current mirror and is subsequently distributed to the conversion stage.

The current-to-code converter generates a multi-bit thermometer output by
replicating $I_2$ and $I_3$ at each comparison node through PMOS and NMOS
current mirrors, respectively, sized with deliberately asymmetric
multiplication factors. In the five-bit implementation described
in~\cite{lee2017}, the PMOS mirror ratios decrease monotonically from
$\times 14$ to $\times 6$ across the five output nodes specifically
$\times 14$, $\times 12$, $\times 10$, $\times 8$, and $\times 6$ for
bits $C_\mathrm{LPT}[4]$ through $C_\mathrm{LPT}[0]$ while the NMOS
mirror ratios increase correspondingly from $\times 6$ to $\times 14$.
At each comparison node, the output logic level is determined by whichever
mirrored current dominates: when the PMOS contribution exceeds the NMOS
contribution, the node is driven to a logic high; otherwise, the NMOS
pull-down prevails and the output settles to a logic low.

To explain the operating principle clearly, it is easier to consider
a configuration with an odd number of output bits, of which the five-bit case
was the chosen one. An odd bit count provides a well-defined, symmetric
midpoint whose comparison threshold aligns exactly with the condition
$V_\mathrm{in} = V_\mathrm{ref}$, making the midpoint of the thermometer
code unambiguous and facilitating a direct account of the circuit behaviour at
the reference voltage. When $V_\mathrm{in}$ lies substantially below
$V_\mathrm{ref}$, the error current $I_\mathrm{err}$ is strongly negative,
meaning $I_2 < I_B$ while $I_3 > I_B$. Since the NMOS mirrors replicate
$I_3$ with larger multiplication factors than the PMOS mirrors replicate $I_2$
at every node, the NMOS current dominates throughout and the output code
resolves to $00000$.

As $V_\mathrm{in}$ rises toward $V_\mathrm{ref}$, the magnitude of
$I_\mathrm{err}$ decreases and the comparison nodes with the highest
PMOS-to-NMOS mirror ratio are the first to transition. The node
$C_\mathrm{LPT}[4]$, carrying a PMOS mirror of ratio $\times 14$ against an
NMOS mirror of only $\times 6$, presents the most favourable balance in favour
of the PMOS side and therefore requires only a modest positive shift in $I_2$
relative to $I_3$ before its output transitions high. Node $C_\mathrm{LPT}[3]$
($\times 12$ PMOS, $\times 8$ NMOS) follows at a moderately larger error
current, while the least significant nodes, $C_\mathrm{LPT}[1]$ and
$C_\mathrm{LPT}[0]$, carry progressively more dominant NMOS mirrors and
require a correspondingly larger current excess before their outputs transition.

At the exact condition $V_\mathrm{in} = V_\mathrm{ref}$, the error current
vanishes and both $I_2$ and $I_3$ are equal to $I_B$. Nodes
$C_\mathrm{LPT}[4]$ and $C_\mathrm{LPT}[3]$ deliver PMOS currents of
$14 I_B$ and $12 I_B$, respectively, against NMOS contributions of $6 I_B$
and $8 I_B$, so both outputs are high. Nodes $C_\mathrm{LPT}[1]$ and
$C_\mathrm{LPT}[0]$ exhibit the reverse relationship and remain low. The
central node $C_\mathrm{LPT}[2]$ carries equal PMOS and NMOS currents of
$10 I_B$ and is therefore in an exactly balanced. The resulting thermometer code
at $V_\mathrm{in} = V_\mathrm{ref,c}$ is consequently $11{?}00$, where the
question mark denotes the indeterminate state of the central bit, fully
consistent with the expected behaviour of a comparator operating precisely at
its decision threshold.

The thresholds governing the transition of each bit are fully determined by
the mirror ratios, as established analytically in~\cite{lee2017}, and result
in a monotonic, symmetric thermometer-coded response across the input voltage
range. This architecture combines inherently fast conversion speed stemming
from the absence of a conventional error amplifier, whose finite bandwidth
would otherwise limit the response time of the control
loop~\cite{lee2017} with a straightforward, hardware-efficient structure and
an explicitly quantifiable input-output characteristic. It was therefore
adopted as the first candidate for the present comparator design, and its
operating principle informed all subsequent topological choices.


\section{Phase II: The Complementary Dual-Input Flash Configuration}

\subsection{Motivation for Wide Input Common-Mode Range Extension}

One of the primary objectives during the development of the proposed comparator architecture was to extend its input common-mode range as close as possible to the supply rails. Achieving a rail-to-rail input capability is desirable because it allows the converter to correctly process analog signals whose common-mode voltage may assume any value within the available supply range, from ground to $V_{DD}$, without requiring additional level-shifting or biasing circuitry.

This characteristic is particularly important in mixed-signal sensor interfaces, where the output voltage of the sensing element is often determined by the physical phenomenon being measured rather than by the requirements of the analog front-end. Restricting the admissible common-mode voltage consequently reduces the effective operating range of the converter and may require additional conditioning stages, increasing both power consumption and circuit complexity.

Furthermore, maximizing the input common-mode range enables better utilization of the available voltage swing. Since the signal occupies a larger fraction of the supply voltage, the effective dynamic range and signal-to-noise ratio can be improved without increasing the supply voltage itself. This consideration becomes especially relevant in modern low-voltage CMOS technologies, where the continuous reduction of $V_{DD}$ inherently limits the available analog headroom.

A wide common-mode operating range also improves the robustness of the system against process, voltage, and temperature (PVT) variations. Small shifts in the input common-mode voltage caused by fabrication mismatch or supply fluctuations are less likely to drive the input transistors out of their intended operating region, thereby preserving the transconductance and ensuring predictable comparator behavior.

For these reasons, a complementary dual-input architecture capable of providing rail-to-rail operation was investigated as a second design candidate.


\subsection{Rail-to-Rail Architecture Evaluation and Voltage Headroom Limitations}

Although the proposed CMF-ADC satisfies the functional requirements of the
event-driven conversion scheme, its input common-mode range remains inherently
limited by the voltage headroom required by the differential input pair and
the tail current source. This limitation becomes particularly critical in
modern low-voltage CMOS technologies, where the continuous reduction of the
supply voltage severely restricts the available analog signal swing. As a
consequence, input signals whose common-mode voltage approaches either supply
rail may drive the transconductance stage outside its intended operating
region, degrading both linearity and sensitivity.

Extending the input common-mode range toward rail-to-rail operation is
therefore highly desirable. Such capability allows the converter to process
analog signals over the entire available supply range without requiring
additional level-shifting or biasing circuitry, increasing compatibility with
a wider variety of sensors and analog front-ends. Furthermore, maximizing the
usable input voltage range improves the utilization of the available dynamic
range, potentially leading to higher signal-to-noise ratio while maintaining
the same supply voltage. It also increases robustness against process,
voltage, and temperature (PVT) variations, since moderate shifts in the input
common-mode voltage are less likely to move the input devices out of
saturation.

Motivated by these considerations, a complementary dual-input architecture was
investigated as a second design candidate. The topology employs two
transconductance stages operating simultaneously: one based on an NMOS
differential input pair and another based on a PMOS differential input pair.
The rationale behind this arrangement is well established in rail-to-rail
input amplifiers~\cite{razavi}. An NMOS differential pair operates efficiently
when the common-mode voltage is close to the negative supply rail, where
sufficient gate overdrive is available to maintain saturation, whereas a PMOS
differential pair naturally provides coverage near the positive supply rail.
By operating both stages simultaneously, at least one transconductor remains
fully functional across the entire supply voltage range.

The implementation explored in this work did not simply combine the output
currents of the two OTAs in parallel. Instead, one output branch of each OTA
was connected to the PMOS current mirrors responsible for generating the
individual bits of the thermometer code, preserving the current-comparison
principle established in the original CMF-ADC architecture. The remaining
output branch of each OTA was then cross-coupled to the complementary
transconductor, such that the residual output current of one stage modulated
the output node of the other, with the same arrangement applied
symmetrically in the opposite direction.

The objective of this bidirectional cross-coupling was to exploit the
complementary transconductance of both stages so that each OTA could reinforce
the operation of the other near the boundaries of its valid input common-mode
range. In principle, this mechanism would extend the effective dynamic
operating region of the converter while preserving the event-driven
current-comparison methodology.

Despite the conceptual advantages of this architecture, its practical
implementation proved incompatible with the target supply voltage of
$V_{\mathrm{DD}} = 0.8~\mathrm{V}$. A conventional five-transistor OTA already
requires a minimum voltage headroom to accommodate the gate-to-source voltage
of the differential pair, the saturation voltage of the tail current source,
and the voltage drop across the active load. At a supply voltage of only
$0.8~\mathrm{V}$, this headroom budget is already severely constrained.

The proposed dynamic-range extension mechanism required the insertion of one
additional transistor in the critical signal path beyond the original
five-transistor topology. Although seemingly a minor modification, every
additional stacked transistor introduces an extra overdrive or saturation
voltage requirement. Consequently, the cumulative voltage headroom exceeded
the available supply, preventing all active devices from remaining in
saturation over the intended operating range and therefore compromising the
correct functionality of the architecture.

An attempt was made to recover the lost headroom by replacing the conventional
current mirrors with the low-voltage current mirror topology proposed by
Razavi~\cite{razavi}. This structure reduces the minimum voltage required
across the mirror branch by eliminating one stacked gate-to-source voltage
drop from the critical path, albeit at the expense of increased circuit
complexity. Nevertheless, even with this modification, the combined voltage
requirements of the differential input pair, the tail current source, the
low-voltage current mirror, and the cross-coupled output stage still exceeded
the available supply voltage.

No feasible bias point could be identified that simultaneously maintained all
transistors within their intended operating regions while preserving the
desired dynamic-range extension mechanism. Consequently, although the proposed
rail-to-rail architecture was conceptually attractive, the stringent headroom
constraints imposed by a $0.8~\mathrm{V}$ supply rendered its practical
implementation unfeasible. For this reason, the architecture was abandoned in
favor of the simpler five-transistor OTA presented in the following section,
which provided a significantly better compromise between functionality,
voltage headroom, implementation complexity, and overall robustness.



\section{Phase III: The Final Adopted Architecture}

\subsection{Five Transistor OTA Core and Optimized Transistor Sizing}

Following the rejection of the complementary dual input approach, the design converged toward a classical five transistor operational transconductance amplifier as the comparator front end. This topology consists of a differential input pair loaded by an active current mirror implemented with two PMOS transistors, where one device operates in diode connection and sets the gate voltage of the mirrored branch. This structure forms the conventional active load of a differential pair and constitutes the defining element of the five transistor OTA architecture~\cite{razavi}.

By eliminating the additional stacked devices required by the previous architecture, this solution satisfies the voltage headroom constraints imposed by the selected supply voltage of $V_\mathrm{DD}=0.8$~V while maintaining correct operation over the intended input range.

The initial implementation evaluated the OTA under open loop conditions in order to maximize gain and improve sensitivity to small differential input voltages. Through transistor sizing optimization, high voltage gain could be achieved and therefore increase comparison sensitivity. However, practical simulations revealed a critical limitation of this approach. Small unintended differential inputs originating from offset, thermal noise, or marginal overdrive conditions caused the output node to rapidly saturate at one of the supply rails before the intended comparison could be completed.

Once saturation occurred, the recovery process became limited by the available bias current responsible for charging or discharging the high impedance output node. This significantly increased settling time and introduced variability in consecutive conversion cycles.

To overcome this limitation, the OTA was reconfigured as a unity gain voltage follower by directly connecting the output node to the inverting input and establishing full negative feedback. Under this configuration, the differential input voltage remains close to zero in steady state and the amplifier continuously operates in its linear region.

As a consequence, the OTA acts as a low impedance replica of the input signal, reproducing $V_\mathrm{in}$ at its output while avoiding saturation effects. The resulting accuracy is determined primarily by the open loop gain and the input referred offset of the amplifier.

\subsection{Ratioed Output Slicers Design}

The quantization stage of the converter is implemented through a bank of CMOS slicer inverters directly connected to the output of the voltage follower. Unlike the previously explored current mirror based architecture, where threshold positioning was determined by current scaling, the present design defines each comparison level through the PMOS to NMOS sizing ratio of each inverter.

The switching threshold of a CMOS inverter is determined by the relative strengths of the pull up and pull down networks and can be approximated by

\begin{equation}
V_M \approx
\frac{V_\mathrm{DD}\sqrt{r}}
{1+\sqrt{r}}
\label{eq:trip_point}
\end{equation}

where

\[
r=\frac{(W/L)_p}{(W/L)_n}
\]

represents the ratio between PMOS and NMOS transistor dimensions.

An increase in $r$ shifts the switching threshold toward higher voltages, whereas lower values move the transition below the mid supply level. By assigning distinct transistor ratios to each inverter, a distributed set of decision thresholds is created along the input voltage axis.

Each slicer therefore acts as an independent comparator operating at a predefined threshold voltage. Collectively, the outputs generate a thermometer coded representation of the input signal.

For the implemented five bit configuration, the most significant output bit, out$_4$, uses a transistor ratio of $r=1.5$, producing a switching threshold above $0.4$~V. In contrast, the least significant bit, out$_0$, adopts $r=0.5$, resulting in a lower threshold. The intermediate slicers are distributed between these limits to generate approximately symmetric threshold spacing around the reference level.

As the input voltage increases from ground toward $V_\mathrm{DD}$, each inverter switches sequentially according to its assigned threshold. At $V_\mathrm{in}=0.4$~V, slicers with $r>1$ remain in logic high state while slicers with $r<1$ have already transitioned to logic low. The inverter designed with $r=1$ operates exactly at its switching point and therefore enters a metastable region.

The resulting thermometer code becomes $11?00$, reproducing the same behavior previously observed in the CMF ADC architecture.

\subsection{Asynchronous Control Logic and Clockless Event Generation Network}

A key characteristic of the proposed architecture is the complete removal of the global clock from the conversion process. Instead of forcing periodic operation, event generation occurs exclusively as a consequence of changes in the analog input signal.

The OTA configured as a voltage follower continuously tracks the input signal and distributes this low impedance representation to all slicers simultaneously. Each slicer independently evaluates whether the input has crossed its corresponding threshold and immediately updates its output state.

The collection of slicer outputs forms an asynchronous thermometer code whose transitions directly represent input activity. Whenever one or more thresholds are crossed, a digital event is naturally generated without requiring any external timing reference.

This event driven operation establishes a direct relationship between signal activity and switching activity inside the converter. During intervals where the input remains constant, the slicers remain static and dynamic power consumption approaches zero. Conversely, periods of rapid input variation automatically increase event generation frequency and conversion activity.

Compared with previous current mirror based implementations, this architecture introduces a clear separation between input buffering and threshold generation. The OTA is responsible only for reproducing the input voltage, while threshold definition is achieved entirely through transistor sizing in the slicer stage.

This separation simplifies the design procedure, improves robustness against saturation, and removes additional voltage headroom constraints, resulting in a compact and scalable architecture suitable for low voltage asynchronous Flash conversion.





\section{Analytical Modeling and Design Constraints}
\subsection{Maximum Input Wave Frequency and Slope Overload}

\subsection{Propagation Delay and Slicer Response Time Modeling}
%confirmar a necessidade desta secção
\subsection{Resolution and Power Dissipation}
%confirmar a necessidade desta secção

\section{Chapter Summary}

